% Options for packages loaded elsewhere
\PassOptionsToPackage{unicode}{hyperref}
\PassOptionsToPackage{hyphens}{url}
\documentclass[
  ignorenonframetext,
]{beamer}
\newif\ifbibliography
\usepackage{pgfpages}
\setbeamertemplate{caption}[numbered]
\setbeamertemplate{caption label separator}{: }
\setbeamercolor{caption name}{fg=normal text.fg}
\beamertemplatenavigationsymbolsempty
% remove section numbering
\setbeamertemplate{part page}{
  \centering
  \begin{beamercolorbox}[sep=16pt,center]{part title}
    \usebeamerfont{part title}\insertpart\par
  \end{beamercolorbox}
}
\setbeamertemplate{section page}{
  \centering
  \begin{beamercolorbox}[sep=12pt,center]{section title}
    \usebeamerfont{section title}\insertsection\par
  \end{beamercolorbox}
}
\setbeamertemplate{subsection page}{
  \centering
  \begin{beamercolorbox}[sep=8pt,center]{subsection title}
    \usebeamerfont{subsection title}\insertsubsection\par
  \end{beamercolorbox}
}
% Prevent slide breaks in the middle of a paragraph
\widowpenalties 1 10000
\raggedbottom
\AtBeginPart{
  \frame{\partpage}
}
\AtBeginSection{
  \ifbibliography
  \else
    \frame{\sectionpage}
  \fi
}
\AtBeginSubsection{
  \frame{\subsectionpage}
}
\usepackage{iftex}
\ifPDFTeX
  \usepackage[T1]{fontenc}
  \usepackage[utf8]{inputenc}
  \usepackage{textcomp} % provide euro and other symbols
\else % if luatex or xetex
  \usepackage{unicode-math} % this also loads fontspec
  \defaultfontfeatures{Scale=MatchLowercase}
  \defaultfontfeatures[\rmfamily]{Ligatures=TeX,Scale=1}
\fi
\usepackage{lmodern}
\ifPDFTeX\else
  % xetex/luatex font selection
\fi
% Use upquote if available, for straight quotes in verbatim environments
\IfFileExists{upquote.sty}{\usepackage{upquote}}{}
\IfFileExists{microtype.sty}{% use microtype if available
  \usepackage[]{microtype}
  \UseMicrotypeSet[protrusion]{basicmath} % disable protrusion for tt fonts
}{}
\makeatletter
\@ifundefined{KOMAClassName}{% if non-KOMA class
  \IfFileExists{parskip.sty}{%
    \usepackage{parskip}
  }{% else
    \setlength{\parindent}{0pt}
    \setlength{\parskip}{6pt plus 2pt minus 1pt}}
}{% if KOMA class
  \KOMAoptions{parskip=half}}
\makeatother
\usepackage{longtable,booktabs,array}
\newcounter{none} % for unnumbered tables
\usepackage{calc} % for calculating minipage widths
\usepackage{caption}
% Make caption package work with longtable
\makeatletter
\def\fnum@table{\tablename~\thetable}
\makeatother
\setlength{\emergencystretch}{3em} % prevent overfull lines
\providecommand{\tightlist}{%
  \setlength{\itemsep}{0pt}\setlength{\parskip}{0pt}}
\usepackage{bookmark}
\IfFileExists{xurl.sty}{\usepackage{xurl}}{} % add URL line breaks if available
\urlstyle{same}
\hypersetup{
  hidelinks,
  pdfcreator={LaTeX via pandoc}}

\author{\texorpdfstring{}{}}
\date{}

\begin{document}

\begin{frame}[fragile]{Methodology: Pipeline Design and Formal
Definitions}
\protect\phantomsection\label{methodology-pipeline-design-and-formal-definitions}
This section describes the five-stage computational meta-analysis
pipeline. Each stage corresponds to a tested, independently executable
script that reads upstream outputs and produces structured artifacts.
The pipeline extends the systematic literature analysis approach of
Knight et al.~\citep{knight2022fep}---which combined manual annotation
with ontology-based automated analysis---by substituting manual coding
with fully automated, LLM-driven assertion extraction and
citation-weighted hypothesis scoring.

\begin{block}{Pipeline Overview}
\protect\phantomsection\label{pipeline-overview}
{\def\LTcaptype{none} % do not increment counter
\begin{longtable}[]{@{}
  >{\raggedright\arraybackslash}p{(\linewidth - 8\tabcolsep) * \real{0.2000}}
  >{\raggedright\arraybackslash}p{(\linewidth - 8\tabcolsep) * \real{0.2000}}
  >{\raggedright\arraybackslash}p{(\linewidth - 8\tabcolsep) * \real{0.2000}}
  >{\raggedright\arraybackslash}p{(\linewidth - 8\tabcolsep) * \real{0.2000}}
  >{\raggedright\arraybackslash}p{(\linewidth - 8\tabcolsep) * \real{0.2000}}@{}}
\toprule\noalign{}
\begin{minipage}[b]{\linewidth}\raggedright
Stage
\end{minipage} & \begin{minipage}[b]{\linewidth}\raggedright
Script
\end{minipage} & \begin{minipage}[b]{\linewidth}\raggedright
Primary Input
\end{minipage} & \begin{minipage}[b]{\linewidth}\raggedright
Primary Output
\end{minipage} & \begin{minipage}[b]{\linewidth}\raggedright
Section
\end{minipage} \\
\midrule\noalign{}
\endhead
1 & \texttt{01\_literature\_search.py} & API queries &
\texttt{corpus.jsonl} & §5.1 \\
2 & \texttt{02\_meta\_analysis\_pipeline.py} & \texttt{corpus.jsonl} &
Classification, temporal, TF-IDF, NMF, citation network JSONs & §5.2 \\
3 & \texttt{03\_build\_knowledge\_graph.py} & \texttt{corpus.jsonl} &
\texttt{nanopublications.jsonl}, \texttt{nanopublications.trig}, scores
& §5.3 \\
4 & \texttt{04\_generate\_figures.py} & All Stage 2--3 JSONs & 16
publication-ready PNGs & §5.4 \\
5 & \texttt{05\_inject\_variables.py} & All output JSONs & Rendered
manuscript Markdown & §5.5 \\
\bottomrule\noalign{}
\end{longtable}
}

Scripts act as thin orchestrators that import methods from tested
library modules and handle file I/O. All computation resides in the
\texttt{src/} packages; no analysis logic is embedded in scripts.
\end{block}

\begin{block}{Stage 1: Multi-Source Literature Retrieval and
Deduplication}
\protect\phantomsection\label{stage-1-multi-source-literature-retrieval-and-deduplication}
We retrieve papers from three complementary academic databases to
maximize coverage and enable cross-source deduplication:

\textbf{arXiv.} We query the arXiv Atom API using phrase-matched
searches including \texttt{all:"active\ inference"},
\texttt{all:"free\ energy\ principle"},
\texttt{all:"expected\ free\ energy"},
\texttt{all:"variational\ free\ energy"\ AND\ all:"inference"}, and
targeted Energy-Based Model queries
(\texttt{all:"energy-based\ model"\ AND\ all:"free\ energy"},
\texttt{all:"Helmholtz\ machine"\ AND\ all:"inference"},
\texttt{all:"Boltzmann\ machine"\ AND\ all:"free\ energy"},
\texttt{all:"contrastive\ divergence"\ AND\ all:"generative\ model"}).
The \texttt{all:} prefix searches titles, abstracts, and full text;
phrase matching reduces contamination from unrelated physics papers that
mention ``free energy'' in thermodynamic contexts. The EBM-adjacent
queries capture research at the intersection of energy-based generative
modeling and variational inference---a growing convergence area
\citep{lecun2006tutorial}.

\textbf{Semantic Scholar.} We query the Semantic Scholar Graph API
\citep{kinney2023semantic} with the same terms. Semantic Scholar
provides citation graphs, abstract embeddings, and links to published
versions. Retry logic with exponential backoff handles rate limiting.

\textbf{OpenAlex.} We query OpenAlex \citep{priem2022openalex} to
capture journal-published work that may not appear on arXiv, including
clinical studies and neuroscience experiments in domain-specific venues.
The \texttt{referenced\_works} field populates citation links for each
paper.

\begin{block}{Canonical Identifier Deduplication}
\protect\phantomsection\label{canonical-identifier-deduplication}
After retrieval, papers are assigned a canonical identifier using the
priority scheme: DOI \(>\) arXiv ID \(>\) Semantic Scholar ID \(>\)
OpenAlex ID \(>\) title hash. When the same paper appears in multiple
sources, the record with the highest metadata completeness is retained.
For each incoming paper, the two records are comparatively evaluated on
metadata completeness---defined as the count of non-empty attributes
among \{abstract, DOI, arXiv ID, venue, citation count\}. The pipeline
reliably retains the structurally richer record; in the event of a tie,
the incumbent is preserved. This ``merge-on-add'' strategy automatically
aggregates the richest available metadata without mandating an expensive
downstream reconciliation pass. Deduplication produces \(N = 785\)
unique papers spanning 1977--2026.
\end{block}

\begin{block}{Relevance Filtering and Curation}
\protect\phantomsection\label{relevance-filtering-and-curation}
After deduplication, a \textbf{relevance filter} removes papers whose
titles and abstracts lack any core Active Inference terminology (e.g.,
\texttt{active\ inference,\textquotesingle{}\textquotesingle{}}free
energy principle,'\,' ``variational free energy'\,'), eliminating
off-topic results introduced by broad keyword overlap across
heterogeneous databases.

We emphasize that this process relies fundamentally on keyword search
strategies across divergent APIs. In any complex research field, there
is no single optimal word or threshold for definitive inclusion or
exclusion. Different information sources and repositories yield
differing schemas and representations, inevitably introducing both false
positives (extraneous papers overlapping in terminology, such as
unrelated database or biological toolkits) and false negatives (relevant
papers employing alternative nomenclature without standard keywords).

Consequently, this pipeline is not intended to produce a static,
``golden'' list of canonical papers. Rather, it is designed as an
open-source software package that can be modularly updated and
versioned. Researchers can configure the pipeline to operate on custom
literature bibliographies curated for specific relevance criteria
through time, treating the initial query-based retrieval as a
programmatic starting point rather than an absolute boundary.
\end{block}
\end{block}

\begin{block}{Stage 2: Bibliometric Analysis}
\protect\phantomsection\label{stage-2-bibliometric-analysis}
Stage 2 performs four complementary analyses on the deduplicated corpus.
All analyses are deterministic given fixed random seeds and operate on
the same \texttt{corpus.jsonl} input.

\begin{block}{Subfield Classification}
\protect\phantomsection\label{subfield-classification}
Each paper is classified into one of eight categories organized across
three domains: \textbf{A -- Core Theory} (A1: quantitative and formal
mathematical theory; A2: qualitative philosophy and general FEP theory),
\textbf{B -- Tools \& Translation} (algorithms, scaling, and software
development), and \textbf{C -- Application Domains} (C1: neuroscience,
C2: robotics, C3: language processing, C4: computational psychiatry, C5:
biology and morphogenesis). Classification uses word-boundary-aware
keyword matching against curated lists applied to titles and abstracts.
A priority system ensures that specific application domains (C1--C5,
priority 1) take precedence over tools (B, priority 2), formal theory
(A1, priority 3), and the broad qualitative philosophy catch-all (A2,
priority 4). Within a priority tier, the domain with the most keyword
matches wins. A1's keyword set includes mathematical indicators such as
\emph{theorem}, \emph{proof}, \emph{convergence}, \emph{posterior},
\emph{equation}, and \emph{Fokker--Planck}, ensuring that papers with
mathematical content are classified as formal theory rather than
defaulting to the philosophy category.
\end{block}

\begin{block}{Temporal Metrics and Growth-Rate Estimation}
\protect\phantomsection\label{temporal-metrics-and-growth-rate-estimation}
We compute temporal publication metrics including year-by-year counts
with gap-filling, cumulative totals, 3-year smoothed moving averages,
and peak year identification. Field dynamics are estimated via two
complementary metrics. The \textbf{mean year-over-year growth rate}
\(\bar{g}\) is the arithmetic mean of annual growth rates for years with
non-zero prior-year publications. The \textbf{doubling time}
\(t_d = \ln 2 / \ln(1 + \bar{g})\). The \textbf{compound annual growth
rate} (CAGR) captures the annualized rate across the full temporal span.
Mathematical details are provided in the Technical Appendix (A.3).
\end{block}

\begin{block}{Text Analytics}
\protect\phantomsection\label{text-analytics}
The TF-IDF matrix is constructed manually using tokenization with
stopword removal and L2-normalized term-frequency
inverse-document-frequency weighting \citep{salton1975vector}, with a
configurable vocabulary size (default: 1000 features). Non-negative
matrix factorization (NMF) is applied to discover latent topics using
multiplicative update rules \citep{lee1999nmf}. Mathematical details are
provided in the Technical Appendix (A.2).
\end{block}

\begin{block}{Citation Network Construction}
\protect\phantomsection\label{citation-network-construction}
The intra-corpus citation network is constructed as a directed graph
where nodes are papers and edges represent citation relationships
resolved within the corpus. Network metrics include PageRank centrality,
HITS hub and authority scores \citep{kleinberg1999authoritative}, degree
distributions, network density, connected components, and community
structure via the Louvain algorithm \citep{blondel2008louvain}.
\end{block}
\end{block}

\begin{block}{Stage 3: Nanopublication-Based Knowledge Graph}
\protect\phantomsection\label{stage-3-nanopublication-based-knowledge-graph}
Stage 3 is the methodological core of this work: it transforms
unstructured abstracts into a structured, RDF-compatible knowledge graph
of scientific evidence. The stage encompasses four tightly coupled
operations: LLM-based assertion extraction, nanopublication packaging,
knowledge graph construction, and citation-weighted hypothesis scoring.

\begin{block}{LLM-Based Assertion Extraction}
\protect\phantomsection\label{llm-based-assertion-extraction}
We extract assertions by prompting a locally hosted LLM (Ollama
\citep{ollama2024}) to assess each paper's abstract against eight
standard hypotheses. The model receives a structured prompt containing
the paper title, abstract, and hypothesis definitions, and returns a
JSON array where each element specifies a hypothesis ID, direction
(supports, contradicts, neutral, or irrelevant), a confidence score
\(c \in [0, 1]\), and a reasoning string. Assertions marked
``irrelevant'' are discarded; confidence values are clamped to
\([0, 1]\); and responses are validated against the known hypothesis ID
set. Papers lacking abstracts are skipped. Detailed prompt engineering,
error taxonomy, and validation methodology are documented in Section 3.
\end{block}

\begin{block}{Nanopublication Schema and RDF Structure}
\protect\phantomsection\label{nanopublication-schema-and-rdf-structure}
Each assertion is encoded as a \textbf{nanopublication}
\citep{groth2010anatomy, kuhn2016decentralized}---a minimal,
self-contained, machine-readable unit of scientific evidence. Formally,
each nanopublication is a tuple \((p, h, d, c)\) where \(p\) is the
paper identifier, \(h\) the hypothesis identifier,
\(d \in \{\text{supports}, \text{contradicts}, \text{neutral}\}\) the
direction, and \(c\) the confidence. Provenance metadata records the LLM
model, UTC timestamp, and paper identifier.

The pipeline serializes nanopublications in two complementary formats:

\begin{enumerate}
\item
  \textbf{JSON Lines} (one JSON object per line) for efficient
  incremental checkpointing. Assertions are saved at configurable
  intervals (default: every 50 papers), enabling the pipeline to resume
  from where it left off after interruption without re-processing
  already-analyzed papers. Deduplication uses the composite key
  \((paper\_id, hypothesis\_id)\); re-runs with improved models
  overwrite stale results.
\item
  \textbf{RDF/TriG} per the nanopublication standard
  (\url{https://nanopub.net/}), producing four named graphs per
  nanopublication:
\end{enumerate}

{\def\LTcaptype{none} % do not increment counter
\begin{longtable}[]{@{}
  >{\raggedright\arraybackslash}p{(\linewidth - 4\tabcolsep) * \real{0.3333}}
  >{\raggedright\arraybackslash}p{(\linewidth - 4\tabcolsep) * \real{0.3333}}
  >{\raggedright\arraybackslash}p{(\linewidth - 4\tabcolsep) * \real{0.3333}}@{}}
\toprule\noalign{}
\begin{minipage}[b]{\linewidth}\raggedright
Named Graph
\end{minipage} & \begin{minipage}[b]{\linewidth}\raggedright
Content
\end{minipage} & \begin{minipage}[b]{\linewidth}\raggedright
Key Predicates
\end{minipage} \\
\midrule\noalign{}
\endhead
\textbf{Head} & Links the nanopub resource to its three component graphs
& \texttt{np:hasAssertion}, \texttt{np:hasProvenance},
\texttt{np:hasPublicationInfo} \\
\textbf{Assertion} & The core scientific claim & \texttt{aif:asserts}
(Paper → Assertion), \texttt{aif:supports}/\texttt{aif:contradicts}
(Assertion → Hypothesis), \texttt{aif:claim}, \texttt{aif:confidence},
\texttt{aif:citationCount} \\
\textbf{Provenance} & How the assertion was generated &
\texttt{prov:wasGeneratedBy}, \texttt{prov:generatedAtTime},
\texttt{prov:wasAttributedTo}, \texttt{prov:hadPrimarySource} \\
\textbf{Publication Info} & Metadata about the nanopublication itself &
\texttt{dc:created}, \texttt{dc:creator}, \texttt{dc:license} \\
\bottomrule\noalign{}
\end{longtable}
}

The namespace \texttt{http://activeinference.org/ontology/} (prefix
\texttt{aif:}) defines all domain predicates; the nanopublication schema
(\texttt{http://www.nanopub.org/nschema\#}, prefix \texttt{np:})
provides structural predicates; provenance uses PROV-O
(\texttt{http://www.w3.org/ns/prov\#}); and Dublin Core
(\texttt{http://purl.org/dc/terms/}) provides publication metadata. The
TriG output is suitable for publication to the decentralized
nanopublication network and aligns with FAIR data principles:
\textbf{F}indable via URI-based identification, \textbf{A}ccessible via
standard RDF protocols, \textbf{I}nteroperable through W3C-standard
serialization, and \textbf{R}eusable with explicit provenance and CC0
licensing.
\end{block}

\begin{block}{Knowledge Graph Construction}
\protect\phantomsection\label{knowledge-graph-construction}
The knowledge graph is an RDF-compatible directed graph with three node
types: \textbf{paper nodes} (metadata: title, abstract, authors, year,
venue, citation count, domain), \textbf{assertion nodes} (claim text,
direction, hypothesis ID, confidence), and \textbf{hypothesis nodes}
(the eight standard hypotheses). Edges encode five relations defined in
the schema:

\begin{itemize}
\tightlist
\item
  \texttt{aif:asserts} --- Paper \(\to\) Assertion
\item
  \texttt{aif:cites} --- Paper \(\to\) Paper
\item
  \texttt{aif:belongsTo} --- Paper \(\to\) Subfield
\item
  \texttt{aif:supports} --- Assertion \(\to\) Hypothesis
\item
  \texttt{aif:contradicts} --- Assertion \(\to\) Hypothesis
\end{itemize}

The graph is implemented with a dual backend: \texttt{rdflib}
\citep{rdflib2023} when available (preferred for semantic web
compatibility), with automatic fallback to \texttt{networkx.DiGraph} for
environments without RDF dependencies. Both backends maintain identical
internal indices for efficient paper, assertion, and hypothesis queries.
\end{block}

\begin{block}{Citation-Weighted Hypothesis Scoring}
\protect\phantomsection\label{citation-weighted-hypothesis-scoring}
For each hypothesis \(H\), we compute a citation-weighted evidence
score:

\[
\text{score}(H) = \frac{\sum_{a \in S(H)} w(a) - \sum_{a \in C(H)} w(a)}{\sum_{a \in A(H)} w(a)}
\]

where \(S(H)\), \(C(H)\), and \(A(H)\) are the sets of supporting,
contradicting, and all assertions for \(H\), and the weight function is:

\[
w(a) = \log(1 + \text{citations}(a)) \cdot \text{confidence}(a)
\]

The logarithmic citation weighting ensures that highly cited papers
carry more influence without allowing any single paper to dominate. The
score lies in \([-1, 1]\). Temporal trends are computed by evaluating
the cumulative score at each year, using only assertions from papers
published up to that year. A full derivation appears in the Technical
Appendix (A.1).
\end{block}

\begin{block}{Tally-Based Evidence Aggregation}
\protect\phantomsection\label{tally-based-evidence-aggregation}
We emphasize that this algorithmic scoring formula constitutes a
\textbf{tally-based approach} to evidence synthesis: each
nanopublication assertion operates as an independent evidential vote,
mathematically weighted by citation impact and the extraction model's
semantic confidence. The aggregation is deliberately linear and
additive---supporting and contradicting assertions are summed and
differenced, independent from modeling dependencies, correlated evidence
clustering, or topological causal structure among claims. This
intentional design choice prioritizes operational transparency, rigorous
reproducibility, and computational tractability over abstract
statistical sophistication.

The tally-based framing introduces three distinct constraints. First,
assertions extracted from methodologically related papers (e.g.,
iterative publications originating from a single research group
validating the same structural model) are counted identically and
independently, inherently amplifying correlated evidence. Second, the
scoring metric imposes symmetrical treatment across assertion source
types: an affirmative assertion parsed from a theoretical review and one
sourced from an empirical randomized controlled trial carry equivalent
leverage at a matched confidence bound. Finally, temporal scoring tracks
\emph{cumulative running totals} rather than dynamic probabilistic
estimates; the score at year \(t\) computes the absolute integrated
momentum of all historical evidence, rather than a decaying posterior
that incrementally downweights early foundational texts.

We embrace these constraints deliberately. The tally-based execution
furnishes a stable, highly interpretable baseline upon which superior
configurations can be systematically evaluated. Section 8 scopes these
concrete extensions---specifically encompassing hierarchical Bayesian
scoring frameworks, causal evidence directed acyclic graphs (DAGs), and
evidential diversity indices that geometrically constrain correlated
research amplification.
\end{block}
\end{block}

\begin{block}{Stage 4: Visualization}
\protect\phantomsection\label{stage-4-visualization}
Stage 4 renders 16 publication-ready figures from the analysis outputs
of Stages 2 and 3. All figures use the Wong (2011) colorblind-safe
palette \citep{wong2011colorblind} and enforce a 16-point minimum font
size for accessibility compliance. Figures span six categories: field
summary and domain distribution (2 figures), growth and temporal
dynamics (2 figures), citation network topology (2 figures), hypothesis
evidence dashboard and timeline (2 figures), assertion composition (2
figures), and text analytics---word cloud, PCA embeddings, term heatmap,
dendrogram, topic-term bars, and co-occurrence matrix (6 figures). The
figure generation script reads only JSON outputs and produces only PNG
files, ensuring strict separation between analysis and visualization.
\end{block}

\begin{block}{Stage 5: Manuscript Variable Injection}
\protect\phantomsection\label{stage-5-manuscript-variable-injection}
Stage 5 computes dynamic variables from all pipeline outputs and injects
them into manuscript Markdown templates via \texttt{\{\{VAR\_NAME\}\}}
placeholder substitution. Variables include corpus-level metrics (size,
year range, CAGR), per-domain counts and percentages, citation network
statistics (nodes, edges, density, components, resolution rate, mean
in-degree), hypothesis scores, and figure counts. All LaTeX-specific
formatting (thousand separators via \texttt{\{,\}}, escaping) is applied
during variable computation, ensuring the manuscript templates remain
human-readable while producing publication-ready LaTeX output.
Unrecognized placeholders are preserved with a warning logged, enabling
incremental manuscript development ahead of full pipeline execution.
\end{block}

\begin{block}{Reproducibility and Test-Driven Validation}
\protect\phantomsection\label{reproducibility-and-test-driven-validation}
The pipeline is deterministic given fixed random seeds and API
responses. Test-driven development enforces 90\% minimum code coverage
on project modules and 60\% on shared infrastructure, with real data and
computation (no mocking). The test suite validates boundary conditions
for hypothesis scoring (all-support \(\to\) +1, all-contradict \(\to\)
\(-1\), balanced \(\to\) 0), schema consistency, serialization
round-trips, and end-to-end pipeline integrity. Source code,
configuration, and outputs are available under CC-BY-4.0.
\end{block}
\end{frame}

\end{document}
